\section{Reduced Dynamics}\label{sec:model}
Consider a free-floating manipulator with a rigid base and $n$ revolute joints, with joint vector $q\in\mathbb R^n$, joint torque $\tau\in\mathbb R^n$, base velocity $\dot x_b=[v_b^\top,\omega_b^\top]^\top\in\mathbb R^6$, and end-effector velocity $\dot x_e\in\mathbb R^6$.

\subsection{Momentum reduction}
With no external force or torque, the kinetic energy is
\begin{equation}
 T=\tfrac12\begin{bmatrix}\dot x_b^\top&\dot q^\top\end{bmatrix}
 \begin{bmatrix}H_b&H_{bm}\\H_{bm}^\top&H_m\end{bmatrix}
 \begin{bmatrix}\dot x_b\\\dot q\end{bmatrix},
 \label{eq:kinetic}
\end{equation}
and the Lagrange equations read
\begin{equation}
 \begin{bmatrix}H_b&H_{bm}\\H_{bm}^\top&H_m\end{bmatrix}
 \begin{bmatrix}\ddot x_b\\\ddot q\end{bmatrix}
 +\begin{bmatrix}C_b&C_{bm}\\C_{bm}^\top&C_m\end{bmatrix}
 \begin{bmatrix}\dot x_b\\\dot q\end{bmatrix}
 =\begin{bmatrix}0\\\tau\end{bmatrix}.
 \label{eq:full}
\end{equation}

\begin{assumption}\label{as:momentum}
The initial linear and angular momenta are zero, i.e.,
$H_b(q)\dot x_b+H_{bm}(q)\dot q=0$, and $H_b(q)$ is uniformly positive
definite on the workspace.
\end{assumption}

Under Assumption~\ref{as:momentum},
\begin{equation}
 \dot x_b=-H_b^{-1}(q)H_{bm}(q)\dot q\triangleq J_{bm}(q)\dot q,
 \label{eq:jbm}
\end{equation}
\begin{equation}
 \dot x_e=\bigl[J_m(q)+J_b(q)J_{bm}(q)\bigr]\dot q\triangleq J_g(q)\dot q,
 \label{eq:gjm}
\end{equation}
where $J_g$ is the generalized Jacobian. Differentiating the momentum
constraint and eliminating $\ddot x_b$ gives the exact equivalent model
\begin{equation}
 M_e(q)\ddot q+C_e(q,\dot q)\dot q=\tau+d(t),
 \label{eq:reduced}
\end{equation}
with
\begin{equation}
\begin{aligned}
 M_e&=H_m-H_{bm}^\top H_b^{-1}H_{bm},\\
 C_e\dot q&=C_m-H_{bm}^\top H_b^{-1}C_b,
\end{aligned}
 \label{eq:schur}
\end{equation}
where $C_b$ and $C_m$ are the bias forces of \eqref{eq:full} evaluated at the
momentum-consistent base velocity \eqref{eq:gjm}. The signal $d(t)$ collects
parameter errors, unmodeled dynamics, and external disturbances.

\begin{assumption}\label{as:inertia}
There is $\underline m>0$ such that $M_e(q)\succeq\underline m I_n$ on the
workspace, and the nominal $M_e$, $C_e$ are available to the controller.
\end{assumption}

\begin{remark}\label{rem:implementation}
Equations \eqref{eq:schur} are implemented exactly: $H_{\mathrm{full}}$ is
assembled with the spatial composite rigid-body algorithm, the bias forces with
a recursive Newton--Euler pass at $\dot x_b=J_{bm}\dot q$, and the Schur
complement and $C_e\dot q$ are extracted afterwards. For the seven-degree-of-freedom
manipulator of Section~\ref{sec:simulation} the implementation returns
$\lambda_{\min}(H_{\mathrm{full}})=0.1375$, a momentum residual
$\|H_b\dot x_b+H_{bm}\dot q\|=4.2\times10^{-15}$, a skew-symmetry residual
$|\dot q^\top(\dot M_e\dot q-2C_e\dot q)|=2.0\times10^{-10}$, and a relative
kinetic-energy drift of $3.3\times10^{-6}$ over one second with $\tau=0$.
\end{remark}

\begin{remark}\label{rem:spread}
For the same manipulator, $\mathrm{diag}(M_e)$ varies between $0.14$ and
$380~\mathrm{kg\,m^2}$, i.e., over a factor of $2.7\times10^3$. This spread,
not the control law, fixes the numerical conditioning of the design in
Section~\ref{sec:control} and the admissible initial error for a given torque
limit.
\end{remark}
